% Page 050

\seite{50}

\abschnitt{Spaziergang}
\pstart
Am 21.\ August unternahmen die vier Schulen der Pfarrei Seffern\ob
gemeinschaftlich einen Spaziergang nach Hamm\unsicher.

\trenner

\pend
\abschnitt{Ernte}
\pstart

Die diesjährige Ernte fiel allgemein befriedigend aus. Die\ob
Grasernte war stellenweise gering. Halmfrüchte wuchsen, auf\ob
die trocknen aber mäßig, weil im Hochsommer lange Zeit\ob
schlechtes Wetter war, vielfach naß eingefahren. Obst gab\ob
es nur in den Tälern, wo die Ernte geringer. Kartoffeln\ob
fielen gut reichlich aus. Die Viehpreise waren hoch.

\trenner

\pend
\abschnitt{Schulfestlichkeiten}
\pstart

Der Geburtstag Sr. Majestät des Kaisers wurde wie alljährlich am\ob
27.\ im Saale zu Seffern gemeinschaftlich von den Schulen der Pfarrei\ob
im Beisein des Herrn Ortsschulinspektors gefeiert. Es wurden\ob
abwechselnd Lieder gesungen und Gedichte vorgetragen. Die\ob
Ansprache an die Kinder hielt Herr Lehrer Winter aus Heilenbach.

\trenner

\pend
\abschnitt{Vaterländische Feier}
\pstart

Am 10.\ März wurde im Pfarrsaale zu Seffern eine gemeinsame Feier der Schü=\ob
len der Pfarrei gehalten, zur Erinnerung an die 100jährige Erhebung Preußens\ob
1813 gegen den Feind. Es wurden abwechselnd Lieder gesungen und Gedichte\ob
vorgetragen. Die Ansprache an die Kinder hielt Lehrer Kremer aus Seffern.

Ges. Seffern 12.\ 1913                  Osterprüfung.\ob
Zimmer\ob
wurde am 9.\ März durch den Ortsschulinspektor Herrn Pfarrer Kimmer aus Seffern\ob
abgehalten.

\pend
\jahresueberschrift{Das Schuljahr 1913-14}
\pstart

Entlassen wurden in diesem Jahre 2 Knaben; dieselben widmen sich der Landwirtschaft. Auf=\ob
genommen wurden 2 Knaben. Die Schülerzahl beträgt 26, davon sind 14 Mädchen, 12 Knaben.\ob
Alle sind katholisch.
\pend
